\section{Discussion and conclusion}
\label{sec:discussion}
\label{sec:conclusion}

We asked what a wider field of view contributes when a lung nodule ranker scores fixed proposals: surrounding anatomy, or more of the lesion. Changing only which native voxels the ranker observes, restoring a thin band around a $10\mm$ ball recovers about 70\% of the CPM gap to the full field for two training seeds and transfers to a second 3D architecture. Sparse access is no substitute: the exact voxels behind the 384 samples close less than half of the gap, and attention over the sample directions adds little. The gains concentrate where the band shows part of the lesion, which happens almost only for lesions of at least $10\mm$ and for proposals away from the lesion centre, as the example in the introduction predicts. Much of the value of wider context at fixed proposals therefore lies in seeing more of the lesion, and context studies should report effects by lesion size and choose patch sizes from the size of the target lesions.

These conclusions have limits. The anatomy analysis is observational and largely coincides with lesion size, and small benefits from beyond the lesion remain compatible with the data. The native-source control changes the number and the contiguity of visible voxels together, and masking creates artificial air boundaries. Three seed-2 attention heads kept their first checkpoint, the native-source and single-stream comparisons use one seed each, and all models share one candidate generator on a reused benchmark. Natural next steps are a band in which the contoured lesion voxels are set to air, an equal-count control that varies only the arrangement of visible voxels, and a test on an untouched cohort with a different proposal stream.
